% !PRINTED 464
\weberpage{464}
%<<CONT
connected with $\overline{P}_2$ on the curve $f$; for it
cannot pass through the points $P_7, P_3$, because in these either
$\Delta$ or $C$ vanishes, but neither can it pass through a
point of the fundamental curve which is not a pole, because it would then
have to pass through the 168 connected points, and could not pass
through the pole $\overline{P}_2$. But then the constant $h$
can be determined so that $K^2 - hQ$ becomes divisible by $f$.

The quotient is an invariant of degree lower than
$42$, but in any case of even degree, and if we therefore
represent it by the fundamental invariants, $K$ cannot
occur in it. From this it follows:

\begin{center}
\gesperrt{The invariant $K^2$ can be expressed rationally in terms of
$f, \Delta, C$.}
\end{center}

If, following this theorem, we represent $K^2$ in the form:
\weq{3}{K^2 = \Sigma\, af^{\nu}\Delta^{\nu_1}C^{\nu_2},}
then in this sum, in which the $a$ are numerical
coefficients, only such terms can occur in which
\[
2\nu + 3\nu_1 + 7\nu_2 = 21,
\]
and by now counting off which values of $\nu, \nu_1, \nu_2$
can occur, we obtain the result that between the
forms
\[
\begin{gathered}
K^2,\quad \Delta^7,\quad C^3,\quad f\Delta^4 C,\quad f^2\Delta C^2,\quad f^3\Delta^5,\\
f^4\Delta^2 C,\quad f^6\Delta^3,\quad f^7 C,\quad f^9\Delta
\end{gathered}
\]
there exists a linear relation with numerical coefficients.
If we represent this relation in the form:
\weq{4}{K^2 = \Phi\Delta + \Psi C,}
in which $\Phi, \Psi$ are likewise invariants, we can draw
from it a further geometric conclusion.

The curves $\Delta$ and $C$ certainly do not intersect on the curve
$f$, because the intersection points of $\Delta$ with $f$ are, all of them, the
$P_7$, those of $C$ with $f$ the $P_3$. But if $\Delta$ and
$C = 0$, then $K = 0$ also, and consequently
\gesperrt{all intersection points of $\Delta$ and $C$ lie on the 21 principal axes of the
group}\footnote{The relation (3) has been completely computed by \gesperrt{Gordan} by the
methods of the theory of invariants (Mathem.\
Annalen, Bd.\ XVII, S.\,371).

We give here the formula there obtained in
the symbols used by us:
\[
K^2 = C^3 - 88f^2\Delta C^2 + 16.63f\Delta^4 C + 17.64f^4\Delta^2 C
- 256f^7 C + 27.64\Delta^7 - 128.469f^3\Delta^5 + 43.512f^6\Delta^3
- 2048f^9\Delta.
\]}.
